\documentclass[11pt,letterpaper]{report}
\usepackage{magent_report}
\usepackage{cleveref}
\newcommand{\KL}{D_{\mathrm{KL}}}
\newcommand{\E}{\mathbb E}
\newcommand{\R}{\mathbb R}
\newcommand{\Prob}{\mathcal P}
\newcommand{\GL}{\mathrm{GL}}
\newcommand{\id}{\mathrm{id}}
\newcommand{\Tr}{\operatorname{tr}}
\newcounter{statementsequence}[chapter]
\newcounter{definition}[chapter]
\newcounter{lemma}[chapter]
\newcounter{proposition}[chapter]
\newcounter{theorem}[chapter]
\newcounter{corollary}[chapter]
\newcounter{remark}[chapter]
\newcounter{openproblem}[chapter]
\newcounter{hypothesis}[chapter]
\makeatletter
\@for\resulttype:=definition,lemma,proposition,theorem,corollary,remark,openproblem,hypothesis\do{%
  \expandafter\edef\csname the\resulttype\endcsname{\noexpand\thechapter.\noexpand\arabic{\resulttype}}}
\makeatother
\newcommand{\resultheading}[4]{%
  \stepcounter{statementsequence}%
  \setcounter{#2}{\numexpr\value{statementsequence}-1\relax}%
  \refstepcounter{#2}\label{#4}%
  \paragraph{#1~\csname the#2\endcsname\ (#3).}}
\newcommand{\definitionheading}[2]{\resultheading{Definition}{definition}{#1}{#2}}
\newcommand{\constructionheading}[2]{\definitionheading{#1}{#2}}
\newcommand{\lemmaheading}[2]{\resultheading{Lemma}{lemma}{#1}{#2}}
\newcommand{\propositionheading}[2]{\resultheading{Proposition}{proposition}{#1}{#2}}
\newcommand{\theoremheading}[2]{\resultheading{Theorem}{theorem}{#1}{#2}}
\newcommand{\corollaryheading}[2]{\resultheading{Corollary}{corollary}{#1}{#2}}
\newcommand{\remarkheading}[2]{\resultheading{Remark}{remark}{#1}{#2}}
\newcommand{\openproblemheading}[2]{\resultheading{Open problem}{openproblem}{#1}{#2}}
\newcommand{\hypothesisheading}[2]{\resultheading{Hypothesis}{hypothesis}{#1}{#2}}
\crefname{definition}{definition}{definitions}
\crefname{lemma}{lemma}{lemmas}
\crefname{proposition}{proposition}{propositions}
\crefname{theorem}{theorem}{theorems}
\crefname{corollary}{corollary}{corollaries}
\crefname{remark}{remark}{remarks}
\crefname{openproblem}{open problem}{open problems}
\crefname{hypothesis}{hypothesis}{hypotheses}
\Crefname{definition}{Definition}{Definitions}
\Crefname{lemma}{Lemma}{Lemmas}
\Crefname{proposition}{Proposition}{Propositions}
\Crefname{theorem}{Theorem}{Theorems}
\Crefname{corollary}{Corollary}{Corollaries}
\Crefname{remark}{Remark}{Remarks}
\Crefname{openproblem}{Open problem}{Open problems}
\Crefname{hypothesis}{Hypothesis}{Hypotheses}
\setcounter{tocdepth}{1}
\makeatletter
\renewcommand{\@pnumwidth}{2.5em}
\renewcommand{\@tocrmarg}{3.5em}
\makeatother

\usepackage{tikz}
\usetikzlibrary{arrows.meta,positioning,calc}
\usepackage{caption}
\captionsetup{font=small,labelfont=bf}
\fancyhead[R]{\small A Physicist's Companion}
\hypersetup{pdftitle={MAgent Theory 2.0: A Physicist's Companion},pdfauthor={Robert C. Dennis},pdfsubject={Physical interpretation, worked examples, and exact constructions for interacting statistical agents}}
\begin{document}
\hypersetup{pageanchor=false}
\magenttitle{MAgent Theory 2.0}{A Physicist's Companion}
\hypersetup{pageanchor=true}
\pagenumbering{roman}
\begingroup
\small\setstretch{1}
\tableofcontents
\endgroup
\bigskip
This companion develops the physical pictures behind statistical agents, their interactions, and inherited organization. It assumes familiarity with probability, linear algebra, and differential equations. Each analogy is tied to a specified mathematical operation. The examples are small enough to calculate by hand; their simplicity does not restrict the general laws to a particular distribution family.

The central problem is how a collective can acquire enough internal organization to continue acting as an agent. Coordinate covariance, a decreasing interaction potential, and a compressed description each answer part of that problem. Their precise roles are developed before considering what would constitute recursive formation or a physical geometry.
\clearpage
\pagenumbering{arabic}
% !TeX root = ../MAgent_Theory_2.0_companion.tex
\chapter{What an Agent Carries}
\label{comp:objects}

A classical statistical description distinguishes a phase-space point from a distribution over phase space. The same distinction begins this theory: $k\in\mathsf K$ is a represented possibility, while $q$ is a probability law over such possibilities. The agent carries the law and mechanisms for changing it. Its contents need not be positions and momenta, and the statistical law need not describe an ensemble of physically separate agents.

\section{Beliefs, models, and predictions}

There are two sectors. The belief law $q_i$ describes state contents $k$. The model law $s_i$ describes model presentations $m\in\mathsf M$. Their most general joint description is
\begin{equation}
\zeta_i(dk,dm)=q_i(dk\mid m)s_i(dm).
\label{comp:eq-joint}
\end{equation}
The notation $q_i(dk\mid m)$ means a conditional law, and $q_i(dk)$ without the condition denotes its marginal. Choosing $\zeta_i=q_i\otimes s_i$ removes dependence between the sectors in that description. It still permits coupled equations for $q_i$ and $s_i$, just as coupled equations for two coordinates do not require a statistical ensemble over those coordinates.

A model presentation becomes predictive through a normalized evaluator $G_i(m,dk)$. For each $m$, this assigns a probability law to the state contents. Averaging the evaluator over the model law gives
\begin{equation}
(s_iG_i)(A)=\int G_i(m,A)s_i(dm).
\label{comp:eq-model-mixture}
\end{equation}
For example, let $k,m$ be binary, with $s_i(m=1)=0.6$, $G_i(1,\{1\})=0.9$, and $G_i(0,\{1\})=0.2$. The model mixture assigns probability $0.62$ to $k=1$. A present belief $q_i(k=1)=0.8$ is a different object. An interaction rule can bring these descriptions into agreement, revise the model law, or change both. Those dynamics must be specified.

An evaluator is a Markov kernel: a function whose output is a probability measure. Ordinary functions are included as the case in which the output measure is concentrated at one value. This language accommodates discrete, continuous, singular, and multimodal laws without selecting a Gaussian family. When differentiating laws later, we will restrict to a regular statistical family or positive finite tables.

\section{A record has an address inside the agent}

Let $x_i\in\mathsf S_i$ denote the complete internal state, including its laws, interfaces, and any memory. A record readout is an ordinary measurable function
\begin{equation}
\rho_i:\mathsf S_i\longrightarrow\mathsf O_i.
\label{comp:eq-readout}
\end{equation}
The output alphabet $\mathsf O_i$ specifies the distinctions the agent can read. A binary register gives $\mathsf O_i=\{0,1\}$; a real-valued comparison can use $\R^r$. The detector analogy is precise at this level: an internal state contains a readable distinction. It does not require an outside observer to inspect it.

Write $x\in\mathsf S$ for the complete population configuration. A deterministic interaction $F_{ij}:\mathsf S\to\mathsf S$ produces the record
\begin{equation}
o_{ij}=\rho_i\bigl(\operatorname{pr}_i F_{ij}(x)\bigr),
\label{comp:eq-deterministic-record}
\end{equation}
where $\operatorname{pr}_i$ selects the receiver's state. For a stochastic transition $\Gamma_{ij}(x,dx')$, the probability of a record in $A$ is
\begin{equation}
\mathcal R_{ij}(x,A)
=\Gamma_{ij}\bigl(x,\{x':\rho_i(\operatorname{pr}_i x')\in A\}\bigr).
\label{comp:eq-record-law}
\end{equation}
This is a pushforward: collect the probability of every next state that displays an outcome in $A$. The transition specifies what happens; the readout specifies what is accessible. A persistent record requires that some later readout still recover the distinction. A transient comparison and a memory trace therefore have different dynamical requirements.

A likelihood on represented contents has another role. For instance, the expression $o=A_i k_i-B_{ij}k_j+\epsilon$ gives a comparison value once the contents and noise are given. Their joint law predicts possible comparisons. A registration rule must still place the outcome in an internal state and specify the response to it.

\section{Which probability space is being used?}

The configuration $x$ contains the agents' laws. A tuple of represented contents $(k_i,m_i)_i$ belongs to a different space. A probability law over possible configurations is a third construction, useful when a preparation or stochastic dynamics has been declared. Deterministic motion of $x$ requires neither sampling the contents nor drawing an ensemble of configurations at each step.

This distinction also clarifies correlations. Two independent fair bits have probability $1/4$ for each pair. Condition on a shared record that the bits agree. Their marginals remain fair, but the conditional joint law has probability $1/2$ on each of $(0,0)$ and $(1,1)$. A calculation of the next joint comparison needs this dependence; the marginals do not contain it. A coupled product-family evolution remains a possible model, but it will not represent that exact conditional joint law. Whether a correlation is needed is decided by the response being predicted.

The laws can be attached to local domains of a context base $\mathcal C$. This is the section picture: a context selects the law applicable there, and a frame specifies its coordinates. Most interactions below occur at a common context. Moving between contexts introduces the connection problem in \cref{comp:geometry}.

% !TeX root = ../MAgent_Theory_2.0_companion.tex
\chapter{Frames, Messages, and Interactions}
\label{comp:links}

An agent's probability law and the coordinates used to describe that law are different objects. Communication introduces a third object: a rule for translating, filtering, or modifying what another agent supplies. Keeping these distinctions explicit prevents a coordinate change from being mistaken for an interaction, or a selective message from being mistaken for an invertible transport.

\section{Reading a link from source to receiver}

There are two sectors. Write $\nu_i^q=q_i$ for the belief marginal and $\nu_i^s=s_i$ for the model marginal. In a finite-dimensional realization, a sector $a\in\{q,s\}$ has a reference vector space $V_a$, an agent space $E_i^a$, and an invertible presentation
\begin{equation}
U_i^a:V_a\longrightarrow E_i^a,\qquad
\bar\nu_i^a=(U_i^a)^{-1}_\#\nu_i^a.
\end{equation}
The pushforward notation expresses the represented possibilities in reference coordinates. The space $E_i^a$ is agent $i$'s displayed copy of $V_a$: at one context all agents' laws lie in one fiber of the statistical bundle, and the presentations are frames used to display them. The two sectors carry two representations of one frame group, with different dimensions, laws, and interaction rules: a belief vector rotates by $\varrho_q(g)$ and a model presentation by $\varrho_s(g)$ under the same change of frame $g$. Think of one set of axes carried by the agent, with beliefs and models as two different kinds of object written on those axes. Each edge can carry a belief link and a model link, both valued in the same group. When they agree the edge transports beliefs and models the same way; when they differ, the mismatch $\tilde\theta_{ij}\theta_{ij}^{-1}$ is a property of the edge that no change of frame removes. The collective geometry of \cref{comp:geometry} uses this single group.

A law-dependent reference generator and its displayed link are
\begin{equation}
L_{ij}^a=\mathcal L_{ij}^a(\bar\nu_i^a,\bar\nu_j^a),\qquad
\Theta_{ij}^a=U_i^a e^{L_{ij}^a}(U_j^a)^{-1}.
\label{eq:comp-link}
\end{equation}
The index order is receiver first, source second. Reading from right to left, the source is expressed in reference coordinates, the generator acts there, and the result is expressed at the receiver. The function $\mathcal L_{ij}^a$ is a constitutive choice. The exponential parameterization does not derive it or equate the belief and model generators. It is also only one kind of interaction. The general interaction is a normalized Markov channel, possibly noisy or information-discarding; the exponential link is the special deterministic channel $y\mapsto\Theta_{ij}^ay$. Its values $e^{L}$ lie in $\exp(\mathfrak{gl}(V_a))$, which in dimension at least two misses some determinant-one matrices, such as $\begin{pmatrix}-1&1\\0&-1\end{pmatrix}$, and every reflection \citep{Culver1966}.

\begin{figure}[htbp]
\centering
\begin{tikzpicture}[
  >={Stealth[length=2.2mm]},
  box/.style={draw=primaryblue,rounded corners=2pt,fill=primaryblue!5,minimum width=2.15cm,minimum height=0.95cm,align=center,font=\small},
  every edge/.style={draw,->,thick,primaryblue},
  every node/.style={font=\small}]
\node[box] (source) {Source fiber\\$E_j^a$};
\node[box,right=1.15cm of source] (refsource) {Reference\\$V_a$};
\node[box,right=1.15cm of refsource] (refreceiver) {Reference\\$V_a$};
\node[box,right=1.15cm of refreceiver] (receiver) {Receiver fiber\\$E_i^a$};
\path (source) edge node[above=3pt] {$(U_j^a)^{-1}$} (refsource)
      (refsource) edge node[above=3pt] {$e^{L_{ij}^a}$} (refreceiver)
      (refreceiver) edge node[above=3pt] {$U_i^a$} (receiver);
\node[below=0.65cm of $(refsource.south)!0.5!(refreceiver.south)$,align=center]
  {$\Theta_{ij}^a=U_i^a e^{L_{ij}^a}(U_j^a)^{-1}$\\[3pt]
   $a=q$ or $s$: two separately specified interactions};
\end{tikzpicture}

\caption{The source presentation is undone by $(U_j^a)^{-1}$, the interaction $e^{L_{ij}^a}$ acts within the reference space, and $U_i^a$ presents the result at the receiver. Endpoint frames and the intervening interaction have separate roles.}
\label{fig:comp-link}
\end{figure}

\section{Changing coordinates}

For a passive change $y_i'=g_i^a y_i$, both the law and its presentation change: $\nu_i^{a\prime}=(g_i^a)_\#\nu_i^a$ and $U_i^{a\prime}=g_i^aU_i^a$. The reference law stays fixed, giving
\begin{equation}
\Theta_{ij}^{a\prime}=g_i^a\Theta_{ij}^a(g_j^a)^{-1}.
\end{equation}
This endpoint covariance is familiar from gauge link variables \citep{Wilson1974}, here in display coordinates. It holds for every generator rule, because the reference laws do not change, and so it constrains nothing: the frames act as a dressing field \citep{FournelEtAl2014}. Every interaction cost can be computed entirely in reference coordinates. Metrics and additional model data must also be re-expressed when forming scalar predictions. Covariance alone supplies no physical symmetry or evolution law.

Changing the common reference presentation is a separate operation. Under $U_i^a\mapsto U_i^a h_a$, the reference laws become $(h_a^{-1})_\#\bar\nu_i^a$. Reference consistency requires the constitutive rule, including any interface data, to give $L_{ij}^{a\prime}=h_a^{-1}L_{ij}^a h_a$. The identity $e^{h_a^{-1}Lh_a}=h_a^{-1}e^Lh_a$ then leaves the displayed link unchanged. This is the substantive symmetry requirement. It is not automatic, and the main text imposes it as an explicit standing hypothesis wherever a conclusion should not depend on the choice of reference frame.

\section{Selective interfaces and an explicit pair}

Sending and receiving can be represented by
\begin{equation}
T_j^a:V_a\longrightarrow H_a,\qquad
O_i^a:H_a\longrightarrow V_a,\qquad
L_{ij}^a=O_i^aT_j^a.
\end{equation}
The sender $T_j^a$ extracts a message; the receiver $O_i^a$ inserts its effect. A direct content map $O_i^aT_j^a$ can lose dimensions. Its exponential is always invertible and acts as identity on $\ker T_j^a$. Thus exponentiating a low-rank generator retains unread directions, and the link compares complete laws. Sending only the extracted message is a different, information-discarding operation. General normalized Markov kernels allow communication without any linear realization.

One explicit law-dependent choice uses finite second moments and positive-definite covariance. Let $\mu_i^a,\Sigma_i^a$ be the reference mean and covariance. On nonself edges choose a scalar $\lambda_{ij}^a$ and set
\begin{equation}
L_{ij}^a=\lambda_{ij}^a\mu_i^a(\mu_j^a)^T(\Sigma_j^a)^{-1}.
\end{equation}
The row $(\mu_j^a)^T(\Sigma_j^a)^{-1}$ sends a scalar message; the column $\lambda_{ij}^a\mu_i^a$ receives it. Under a reference change, the mean transforms by $h_a^{-1}$ and covariance by congruence. Their product consequently transforms by conjugation. These assumptions concern moments, not Gaussian shape. For these links the primitive self link is set to the identity separately. In general the self channel is constitutive data, and some later constructions use a nonidentity one.

For a concrete non-Gaussian example, let $Z$ have independent centered Laplace coordinates of variance one in two dimensions. Give two agents the laws of $e_1+Z$ and $e_2+Z$, with identity frames and unit coefficients. Then $\Sigma_1=\Sigma_2=I$, $L_{12}=E_{12}$, and $L_{21}=E_{21}$. Since each generator squares to zero,
\begin{equation}
\Theta_{12}\Theta_{21}
=(I+E_{12})(I+E_{21})
=\begin{pmatrix}2&1\\1&1\end{pmatrix}.
\end{equation}
Here $E_{ab}$ has its only nonzero entry, one, in row $a$, column $b$. The directed communication loop is nontrivial. It establishes neither a connection over contexts nor a physical curvature. Writing $\Theta_{12}=\Omega_{12}N_2^{1/2}$ with $N_2=\Theta_{21}\Theta_{12}$, the reciprocal part $\Omega_{12}$ is here a rotation with $\Omega_{12}\Omega_{21}=I$, so this two-cycle measures nonreciprocity, not holonomy.

A further limit concerns what any cost of the form $\KL(\nu_i\Vert\nu_j\mathcal K_{ij})$ can see. A loop map $H$ enters only through the transported law $H_\#\nu$. If $H$ fixes the law, as a rotation fixes a centered isotropic Gaussian, the cost cannot register it. The joint law of a single sample $Y$ and its image $HY$ does determine $H$ almost everywhere, so detecting such a loop requires a record that compares one sample along two routes.

\chapter{Coupled Relaxation and Source Selection}
\label{comp:dynamics}

The dynamical variables are probability laws. Moving an agent therefore means changing its expectations, uncertainties, and dependencies, rather than moving a point through its sample space. The physical analogue developed here is overdamped motion on a probability manifold: an interaction potential supplies the driving covector, and the Fisher metric converts that covector into a velocity. Neither ingredient determines the other.

\section{The potential and its source rows}

At a fixed context, let $q_i$ describe agent $i$'s represented state $k$, and let $s_i$ describe its model presentation $m$. For this construction their joint recognition law is the explicit product $q_i\otimes s_i$. Choose a model reference $r_i(dm)$ and a normalized evaluator $G_i(m,dk)$. Their local contribution is
\begin{equation}
 \mathcal E_0=\sum_i\KL(q_i\otimes s_i\Vert r_iG_i)
 =\sum_i\left[\KL(s_i\Vert r_i)
       +\int\KL(q_i\Vert G_i(m,\cdot))s_i(dm)\right].
 \label{eq:comp-local-potential}
\end{equation}
The equality separates the model marginal from the conditional state law. Thus a product recognition family still has coupled dynamics: changing $s_i$ changes the evaluator against which $q_i$ relaxes, and changing $q_i$ changes the model costs.

The channels of \cref{comp:links} send peer laws into comparable spaces. Write $\nu^q_{ij}=q_j\mathcal K^q_{ij}[\mathbf q]$ and $D^q_{ij}=\KL(q_i\Vert\nu^q_{ij})$, where $\mathbf q$ collects the belief laws. Define $\nu^s_{ij}$ and $D^s_{ij}$ similarly for model laws. With normalized source rows $\beta_i,\gamma_i$, positive reference rows $\pi_i^q,\pi_i^s$, and $T_q,T_s>0$, the potential is
\begin{equation}
 \begin{aligned}
 \mathcal E=\mathcal E_0
 &+\sum_i\left[\sum_j\beta_{ij}D^q_{ij}
                  +T_q\KL(\beta_i\Vert\pi_i^q)\right]\\
 &+\sum_i\left[\sum_j\gamma_{ij}D^s_{ij}
                  +T_s\KL(\gamma_i\Vert\pi_i^s)\right].
 \end{aligned}
 \label{eq:comp-potential}
\end{equation}
Each sum uses its declared source set. The row penalty balances preference for compatible peers against departure from the reference routing.

Hold the laws fixed. For either sector, minimize $f(\eta)=\sum_j\eta_jD_j+T\KL(\eta\Vert\pi)$ subject to $\sum_j\eta_j=1$. Stationarity gives $D_j+T[\log(\eta_j/\pi_j)+1]=\lambda$, hence
\begin{equation}
 \eta_j^*=\frac{\pi_je^{-D_j/T}}{Z},\qquad
 Z=\sum_j\pi_je^{-D_j/T},\qquad
 f(\eta)=-T\log Z+T\KL(\eta\Vert\eta^*).
 \label{eq:comp-source-row}
\end{equation}
The last identity proves that this stationary row is the unique minimum. Eliminating both sets of rows gives $\mathcal E_{\mathrm{red}}=\mathcal E_0-T_q\sum_i\log Z_i^q-T_s\sum_i\log Z_i^s$. At fixed temperatures and reference rows,
\begin{equation}
 d\mathcal E_{\mathrm{red}}
 =d\mathcal E_0+\sum_{ij}\beta_{ij}^*dD^q_{ij}
                   +\sum_{ij}\gamma_{ij}^*dD^s_{ij}.
 \label{eq:comp-envelope}
\end{equation}
This follows by differentiating each log partition sum. Optimized row derivatives disappear; derivatives inside their costs remain.

\section{Every interaction role contributes to motion}

For normalized densities $a,b$, variation of $\int a\log(a/b)$ gives
\begin{equation}
 d\KL(a\Vert b)=\int\log(a/b)da-\int(a/b)db,
 \qquad
 d\nu^q_{ij}=(dq_j)\mathcal K^q_{ij}
             +q_jD\mathcal K^q_{ij}(d\mathbf q).
 \label{eq:comp-all-role-variation}
\end{equation}
Here $da,db$ denote mass-preserving variations, and differentiation assumes a regular dominated family. An agent moves through its receiver role, its source role, and any dependence of the channels on its law. Directed routing does not remove incoming contributions to source motion. For an identity self-channel, the two laws remain identical and their divergence stays zero. Holding an old self snapshot fixed instead creates a different potential after the current law changes.

For positive finite probability tables, collect all dynamical tables as $u^a$, including the source rows when they are not eliminated. The Fisher metric and its unit-mobility descent are
\begin{equation}
 \begin{aligned}
 g_u(v,v)&=\sum_{a,z}\frac{(v_z^a)^2}{u_z^a},\\
 \dot u_z^a&=-u_z^a(h_z^a-\E_{u^a}h^a),\qquad
 h_z^a=\frac{\partial\mathcal E}{\partial u_z^a},\\
 \dot{\mathcal E}&=-\sum_a\operatorname{Var}_{u^a}(h^a).
 \end{aligned}
 \label{eq:comp-fisher-descent}
\end{equation}
Centering preserves each normalization. Substituting the velocity into $\dot{\mathcal E}=\sum_{a,z}h_z^a\dot u_z^a$ proves the dissipation identity. The same construction uses the Fisher tensor on a regular smooth family \citep{Amari1998}. Both sectors have equal mobility; a slower model response must arise from the resulting dynamics.

A Bernoulli law makes the metric concrete. Its success probability $p$ has metric $dp^2/[p(1-p)]$, so the mobility is $p(1-p)$. For $\mathcal E(p)=\KL(\operatorname{Ber}(p)\Vert\operatorname{Ber}(p_0))$, with $p,p_0\in(0,1)$, write $l=\log[p/(1-p)]$ and $l_0=\log[p_0/(1-p_0)]$. Then
\begin{equation}
 \dot p=-p(1-p)(l-l_0),\qquad
 \dot l=-(l-l_0).
 \label{eq:comp-bernoulli-flow}
\end{equation}
Log odds relax exponentially, although probability does not follow a linear restoring equation. Here $t$ parametrizes the chosen relaxation, and the dimensionless $T_q,T_s$ weight routing penalties. They have not been identified with physical time or bath temperature.

\section{A record likelihood defines a different objective}

A fixed normalized prior $P_0$ and record likelihood $\ell_o$ give posterior $\Pi_o=\ell_oP_0/Z_o$ and the exact identity
\begin{equation}
 \mathcal F(Q;o)=\KL(Q\Vert P_0)-\E_Q\log\ell_o
               =\KL(Q\Vert\Pi_o)-\log Z_o.
 \label{eq:comp-record-vfe}
\end{equation}
Source inference in a product recognition family then uses expected negative log likelihood as its cost. Peer coupling instead uses divergence between laws, whose source template can move during relaxation. Both yield Gibbs rows because both contain the same categorical relative-entropy penalty; their costs and state derivatives differ. Identical independent Gaussian laws have zero peer divergence but nonzero expected squared sample difference. A normalized likelihood for that difference therefore describes an interaction distinct from direct agreement of the laws.

% !TeX root = ../MAgent_Theory_2.0_companion.tex
\chapter{What a Message Lets the Receiver Infer}
\label{comp:communication}

Sending a few observables can change an agent's description of many variables. The extra structure comes from the receiver's conditional model. This is familiar in physics: measuring one coordinate of a correlated system can change predictions for another coordinate without measuring it directly. The distinction between transmitted information and inferred structure is explicit here.

\section{Conditional prediction within an agent}
\label{comp:conditional-prediction}

An agent can use its model to infer a result from a different, correlated cue. Let $A\in\{1,\ldots,M\}$ be a finite model presentation with prior weights $\pi_a>0$ summing to one. On standard Borel cue and result spaces, presentation $a$ supplies a cue law $\mu_a(du)$ and a normalized result kernel $K_a(u,dv)$. A normalized record kernel $\mathcal L(u,do)$ measures the cue. The joint law is
\begin{equation}
P(a,du,dv,do)=\pi_a\mu_a(du)K_a(u,dv)\mathcal L(u,do).
\label{comp:eq-cue-result-joint}
\end{equation}
Successive integration of the kernels and summation over $a$ give total mass one. In the agent description, $a$ is the model presentation, $(u,v)$ is the represented content, and $\mu_a(du)K_a(u,dv)$ is its evaluator. The kernels and prior weights constitute the stored model; the posterior after a record describes the agent's current beliefs.

Suppose $\mathcal L(u,do)=l(o\mid u)\lambda(do)$ and the observed $o$ has finite positive evidence $Z_a(o)=\int l(o\mid u)\mu_a(du)$ in every component. Bayes' rule gives
\begin{equation}
\begin{aligned}
w_a&=\frac{\pi_aZ_a(o)}{\sum_b\pi_bZ_b(o)},\\
\Pi_a^o(du,dv)&=\frac{l(o\mid u)\mu_a(du)K_a(u,dv)}{Z_a(o)}.
\end{aligned}
\label{comp:eq-cue-result-posterior}
\end{equation}
The exact recognition law is $w_a\Pi_a^o$, so the model belief $s(a)=w_a$ is generally correlated with the contents. A product restriction would replace the component-dependent content laws by one common law. Conditioning revises these beliefs without changing the stored kernels. This construction uses general probability laws; Gaussian conditioning makes a particular realization explicit.

For a scalar example, choose one presentation and independent variables
\begin{equation}
X\sim\mathcal N(0,1),\qquad
\epsilon\sim\mathcal N(0,1/4),\qquad
\eta\sim\mathcal N(0,1),
\end{equation}
with result $Y=-2X+\epsilon$ and record $O=X+\eta$. Observe $O=\xi$. The cue posterior follows by completing the square:
\begin{equation}
p(x\mid\xi)\propto
\exp\left[-\frac{x^2+(\xi-x)^2}{2}\right]
\propto\exp[-(x-\xi/2)^2],
\end{equation}
hence $X\mid O=\xi\sim\mathcal N(\xi/2,1/2)$. Since the result noise is independent of the cue and record,
\begin{equation}
Y\mid O=\xi\sim\mathcal N(-\xi,9/4).
\label{comp:eq-negative-result}
\end{equation}
The variance is $4(1/2)+1/4=9/4$. Increasing the observed cue increases the inferred cue but decreases the predicted result. The negative response comes from the model's conditional relation between two variables. Estimating the cue itself in this example would instead have positive slope $1/2$. Conditional prediction therefore gives a statistical account of a transformation that reverses a direction. Conditioning a joint Gaussian mixture to predict one group of variables from another is Gaussian mixture regression \citep{Calinon2016}.

With several presentations, uncertainty also concerns which relation applies. For vector-valued results with finite second moments, let $m_a=\E[Y\mid o,A=a]$ and $S_a=\operatorname{Cov}(Y\mid o,A=a)$. The result marginal $p_Y^o$ has moments
\begin{equation}
\begin{aligned}
\bar m&=\sum_a w_am_a,\\
S&=\sum_a w_a\left[S_a+(m_a-\bar m)(m_a-\bar m)^\top\right].
\end{aligned}
\label{comp:eq-result-ambiguity}
\end{equation}
The first covariance contribution is uncertainty within a presentation. The second is uncertainty from the disagreement of its predictions with those of other presentations. These terms follow by writing $Y-\bar m=(Y-m_a)+(m_a-\bar m)$ and taking the conditional expectation; the cross terms vanish.

If the result is a finite mixture of nonsingular Gaussians, one possible carrier is $q_G=\mathcal N(\bar m,S)$. It minimizes $\KL(p_Y^o\Vert q_G)$ over Gaussian laws. Indeed, the expected negative Gaussian log density depends on $p_Y^o$ only through its first two moments, and its minimum occurs at their matching mean and covariance. This is the moment projection used in assumed-density filtering \citep{Minka2001EP}. The variational free energy instead contains the divergence with recognition in the first slot. Minimizing $\KL(q\Vert p_Y^o)$ over Gaussians is a different problem and need not return the moment carrier.

Retaining $w$ alongside $q_G$ still does not preserve their original dependence. If the retained description is the product $w\otimes q_G$, the exact label/result posterior $\Pi_{A,Y}^o$ obeys
\begin{equation}
\KL(\Pi_{A,Y}^o\Vert w\otimes q_G)
=I_{\Pi^o}(A;Y)+\KL(p_Y^o\Vert q_G).
\label{comp:eq-memory-retention-loss}
\end{equation}
Here $I_{\Pi^o}(A;Y)=\KL(\Pi_{A,Y}^o\Vert w\otimes p_Y^o)$ is mutual information. Insert $w\otimes p_Y^o$ into the logarithmic density ratio to obtain the identity. It separates the loss of model/content dependence from the loss of non-Gaussian result information. Neither loss is determined by the accuracy of the mean alone.

A later step may carry the component belief forward, for example as $wT$ under a declared stochastic transition matrix $T$. An interpretation as Bayesian filtering also needs the conditional law of the next record given the retained history and presentation. A deterministic function of a record already available to the agent supplies no further evidence conditional on that record; treating it as an independent observation would change the model. Communication can supply information new to the receiver, but its record kernel must be specified. The local conditioning construction therefore provides an agent computation, while repeated assimilation, model learning, and population dynamics require their own declared mechanisms.

\section{Selective sending and conditional reconstruction}

A sender uses a measurable map $T_j:E_j\to\mathsf H$ into a message space. If its content has law $q_j$, possible messages have law
\begin{equation}
\mu_j=(T_j)_\#q_j.
\label{comp:eq-message-law}
\end{equation}
The pushforward symbol means $\mu_j(A)=q_j(T_j^{-1}A)$. Sampling a message and transmitting the whole law $\mu_j$ are different communication protocols. Messages run from sender to receiver. The reciprocal potential of the dynamics chapter also moves each sender toward its listeners, which a communicating population can only realize if the listeners send something back. A linear receiving map $O_i:\mathsf H\to E_i$ embeds a message, but its composition with $T_j$ cannot have higher rank than $T_j$.

\begin{figure}[htbp]
\centering
\begin{tikzpicture}[
  >={Stealth[length=2.2mm]},
  box/.style={draw=primaryblue,rounded corners=2pt,fill=primaryblue!5,minimum width=2.55cm,minimum height=1.05cm,align=center,font=\small},
  every node/.style={font=\small}]
\node[box] (source) {Sender law\\$q(dk)$};
\node[box,right=1.6cm of source] (message) {Message law\\$\nu=T_\#q$};
\node[box,right=1.6cm of message] (receiver) {Reconstruction\\$\widehat q(dk)$};
\node[box,below=1.1cm of receiver,fill=secondaryblue!10] (prior) {Receiver conditional\\$p(dk\mid h)$};
\draw[->,thick,primaryblue] (source) -- node[above=4pt] {$T$} (message);
\draw[->,thick,primaryblue] (message) -- node[above=4pt,align=center] {mix over $h$} (receiver);
\draw[->,thick,secondaryblue] (prior) -- (receiver);
\node[below=0.65cm of source,align=center,text width=3cm] {Different sender laws\\can give the same $\nu$.};
\end{tikzpicture}

\caption{Selective communication and receiver structure. The sender fixes the transmitted marginal. The receiver supplies a conditional law for distinctions that the message does not specify. This distributional reconstruction presupposes that the message law is available; one received sample supplies a value instead.}
\label{fig:comp-message}
\end{figure}

To see exactly what is lost, work in one common presentation and write $T:E\to\mathsf H$. Let $q$ be the sender's law and $p$ the receiver's reference law, with finite $\KL(q\Vert p)$. The receiver can combine the transmitted marginal $\mu=T_\#q$ with its own conditional law:
\begin{equation}
\widehat q(dk)=\int p(dk\mid h)\mu(dh).
\label{comp:eq-reconstruction}
\end{equation}
Disintegration on standard Borel spaces provides these conditional laws \citep{Kallenberg2021}. The reconstruction has exactly the right message marginal. Its remaining discrepancy is
\begin{equation}
\KL(q\Vert\widehat q)
=\int\KL\bigl(q(\cdot\mid h)\Vert p(\cdot\mid h)\bigr)\mu(dh).
\label{comp:eq-reconstruction-discrepancy}
\end{equation}
Both laws assign the same probability to $h$, so that factor cancels in their density ratio. The remaining logarithm compares the two conditional laws, which gives the equality by integration. This reconstruction is Jeffrey's rule of conditioning \citep{DiaconisZabell1982}, and it is the law closest to the receiver's prior among all laws with the transmitted marginal. The divergence then splits like a right triangle, $\KL(q\Vert p)=\KL(q\Vert\widehat q)+\KL(\widehat q\Vert p)$, with the second leg equal to the divergence of the message marginals \citep{Csiszar1975}.

For a calculation on four points, let $k=(H,V)$ be two bits. The sender gives probability $1/2$ to each of $00$ and $11$, and transmits only $H$. A receiver whose reference makes both bits independent and fair reconstructs the uniform law on all four points. The message marginal is correct, while the relative entropy is $\log2$. A receiver whose conditional model has $V=H$ reconstructs the sender's full law. The missing correlation was supplied by the receiver.

One reference conditional can recover a whole family only under a stronger condition. Suppose $q_\alpha=L_\alpha p$ and use the same $p(dk\mid h)$ for every $\alpha$. Its reconstruction has density
\begin{equation}
\frac{d\widehat q_\alpha}{dp}(k)
=\E_p[L_\alpha\mid T](k).
\end{equation}
To obtain this identity, weight the reference message marginal by its conditional mean likelihood and then apply the reference conditional. The result equals $q_\alpha$ precisely when $L_\alpha$ is already measurable from the message, which is the classical statement that the message is a sufficient statistic for the family \citep{HalmosSavage1949}. Reconstructing one posterior with its own conditional does not establish this family-wide property. If a local reconstruction must also ignore retained exterior information, the reference needs conditional independence from that exterior given the message.

For a single received value $h$, Bayesian assimilation instead takes the form
\begin{equation}
q_i^+(dk)=\frac{\ell_i(h\mid k)q_i^-(dk)}{\int\ell_i(h\mid k')q_i^-(dk')},
\label{comp:eq-bayes-receiver}
\end{equation}
with a positive finite denominator. In the perfectly correlated two-bit example, a noiseless observation $H=1$ also fixes $V=1$. This is inference through an existing relation, not an inversion of an arbitrary lossy map. Under the reciprocal potential such an update is a kick: the flow afterwards relaxes toward states that do not depend on the record, unless the record's likelihood is kept in the objective.

\section{Agreement needs a common prediction rule}

Suppose a model law $s$ and a normalized record evaluator $K(m,do)$ predict the record law $sK$. If two agents compare model laws in one frame, they must also compare their evaluators in that frame. Equal model laws can predict opposite bits when their evaluators differ.

For one common evaluator and finite $\KL(s\Vert t)$, an exact information balance is
\begin{equation}
\KL(s\Vert t)
=\KL(sK\Vert tK)
+\int\KL(s^o\Vert t^o)(sK)(do),
\label{comp:eq-prediction-balance}
\end{equation}
where $s^o,t^o$ are the conditional model laws after record $o$. The joint laws $s(dm)K(m,do)$ and $t(dm)K(m,do)$ have divergence $\KL(s\Vert t)$ because their common kernel cancels. Factoring the same joint laws by their record marginal gives the displayed split. The last term is the distinction left unresolved by the record. A constant evaluator discards every distinction; a sufficient record leaves none.

For a model link $\Theta^s_{ij}$, compatibility means $K_i(\Theta^s_{ij}m,A)=K_j(m,A)$ after identifying record alphabets. Only then does transporting the model law also transport its predictive meaning. If the link changes with the laws, this condition must continue to hold at the updated state.

\section{Compression should name what matters}

Rate--distortion theory makes the desired fidelity part of the communication problem \citep{Shannon1959}. A low-dimensional real vector alone is not a finite-bit channel. Noise, precision, or an encoding constraint must specify the communication cost. Wyner--Ziv coding additionally gives the receiver a realized variable correlated with the source \citep{WynerZiv1976}; possession of a prior by itself is not that realized side information.

For a predictive criterion, name the source $X$, target response $Y$, and accessible receiver information $Z$. An encoder generates $U$ from $X$ alone, so $U-X-(Y,Z)$ is a conditional-independence chain. The optimal decoder is $P(Y\mid U,Z)$. Its expected relative-entropy distortion is
\begin{equation}
\mathcal D=I(Y;X\mid U,Z)
=I(Y;X\mid Z)-I(Y;U\mid Z).
\label{comp:eq-predictive-distortion}
\end{equation}
The equality follows by expanding $I(Y;X,U\mid Z)$ in both orders. It is the information about the chosen response that compression discards. An information-bottleneck criterion trades this quantity against a cost such as $I(X;U\mid Z)$ \citep{Tishby1999}; an operational bit rate requires additional coding assumptions.

If $X=(A,B)$ consists of independent fair bits and $U=A$, the distortion is zero for target $Y=A$ and $\log2$ for target $Y=B$. For a prospective parent, the target can be a future boundary response. Retaining one such response does not ensure that every later interaction can be inherited. That stronger requirement concerns dynamics.

\chapter{When a Collective Can Act as an Agent}
\label{comp:inheritance}

A collective can have a simple description without being able to evolve from
that description alone. Knowing the total energy of two bodies, for example,
does not usually determine their next interaction. A parent of law-valued
agents faces the same question: does its retained state determine its future
responses?

There are two operations to distinguish. A channel $C$ compresses represented
contents and gives a parent law $Q_xC$. A map $R$ retains a complete dynamical
state $R(x)$. Besides belief and model laws, that state can include the model
evaluator, source information, boundary channels, and memory. A distribution
over possible contents does not automatically contain these additional data.

\section{Two routes to the same parent}

Suppose the fine configuration advances by $F$. Exact inherited evolution
requires a parent update $F^c$ such that
\begin{equation}
 R(F(x))=F^c(R(x)).
\end{equation}
We can first evolve the children and then form their parent, or first form
the parent and evolve it. Both routes must give the same retained state.
Repeating the identity gives agreement at every later step. Boundary response
mechanisms must likewise depend only on the retained state.

\begin{figure}[t]
\centering
\begin{tikzpicture}[
  >={Stealth[length=2.2mm]},
  box/.style={draw=primaryblue,rounded corners=2pt,fill=primaryblue!5,minimum width=2.55cm,minimum height=0.95cm,align=center,font=\small},
  every node/.style={font=\small}]
\node[box] (fine) {Fine state\\$x$};
\node[box,right=4.3cm of fine] (finenext) {Fine successor\\$F(x)$};
\node[box,below=1.8cm of fine,fill=secondaryblue!10] (parent) {Parent state\\$R(x)$};
\node[box,below=1.8cm of finenext,fill=secondaryblue!10] (parentnext) {Parent successor\\$R(F(x))=F^c(R(x))$};
\draw[->,thick,primaryblue] (fine) -- node[above=4pt] {fine update $F$} (finenext);
\draw[->,thick,secondaryblue] (parent) -- node[below=4pt] {inherited update $F^c$} (parentnext);
\draw[->,thick,primaryblue] (fine) -- node[left=5pt] {$R$} (parent);
\draw[->,thick,primaryblue] (finenext) -- node[right=5pt] {$R$} (parentnext);
\node at ($(fine)!0.5!(parentnext)$) {same retained result};
\end{tikzpicture}

\caption{The two routes to the next parent configuration must agree.
Vertical maps retain information; horizontal maps advance the dynamics.
A compressed law alone need not make this square commute.}
\label{fig:comp-inheritance}
\end{figure}

The immediate test is whether configurations with the same parent have the
same next parent. If $R(x)=R(x')$ but $R(F(x))\ne R(F(x'))$, no single
$F^c$ can satisfy the equation. Conversely, in a finite state space, this
test is sufficient: evolve any representative of each parent state.
Agreement within each fiber of $R$ makes the choice irrelevant. On general
measurable spaces the resulting map must also be measurable.

For stochastic dynamics, replace the next state by its distribution. If $T$
is the fine transition matrix, the probability
\begin{equation}
 \sum_{y:R(y)=z'}T(x,y)
\end{equation}
must depend only on $R(x)$. These sums then define normalized parent rows.
Conditioning on any parent history merely averages identical rows, so the
parent reproduces the whole retained history law. This is strong lumpability
\citep{KemenySnell1976}. Agreement for one equilibrium preparation is a
weaker statement.

\section{When a gradient flow descends}

Suppose the fine configuration lies on a Riemannian manifold with metric $g$ and evolves by $-\operatorname{grad}_g E$. Let a smooth surjective submersion $R$ retain the parent, whose metric is $g_c$. One sufficient condition is that the energy depend only on the parent,
\begin{equation}
E=E_c\circ R,
\end{equation}
and that $R$ preserve lengths of horizontal vectors, meaning vectors orthogonal to $\ker DR$. Then the gradient is horizontal because the energy has no derivative along a fiber of $R$. For a horizontal $u$,
\begin{equation}
g_c(DR\operatorname{grad}_gE,DR u)
=g(\operatorname{grad}_gE,u)
=dE_c(DR u).
\end{equation}
Every parent tangent has a horizontal lift, so the pushed gradient is exactly $\operatorname{grad}_{g_c}E_c$. The dynamics descends in the same time.

If instead $E=\chi(E_c\circ R)$ and the horizontal metric is scaled by $d(x)^2$, the pushed vector field acquires the factor $\chi'(E_c(Rx))d(x)^2$. A positive factor preserves oriented trajectories after a time change. A factor that depends on discarded state does not specify an autonomous parent clock. Taking the minimum of $E$ over each fiber does not supply the first hypothesis: it gives the value at the minimizing states, not the fine energy everywhere on that fiber.

The requirement $E=E_c\circ R$ is stronger than descent needs, and it has a physical cost: it makes the energy blind to every motion along a fiber, so the discarded coordinates cannot move relative to the retained ones. Only the horizontal derivatives matter. The computation above uses $dE(u)$ for horizontal $u$ alone, so descent holds exactly when
\begin{equation}
dE(u)=dE_c(DR(u))\qquad\text{for every horizontal }u.
\end{equation}
Two particles on a line with unit masses, harmonic self-energies $\tfrac12x_1^2+\tfrac12x_2^2$, and any pair interaction $W(x_1-x_2)$ show the difference. The center of mass $(x_1+x_2)/\sqrt2$ relaxes as if the interaction were absent, because a force that depends only on the separation pushes the two particles in opposite directions and cancels in their sum. Meanwhile the separation evolves under $W$. The parent is exact, the discarded coordinate has its own dynamics, and the energy is not a function of the center of mass, except for the single interaction $W(s)=-s^2/4$ (up to a constant) that exactly cancels the separation part of the harmonic self-energies and freezes the separation.

The opposite case is a retention that simply drops some agents. If the retained and discarded agents exchange any peer cost, the retained velocity depends on the discarded state, because each agent pulls on every agent it attends to and is pulled back. Dropping whole agents is exact only for blocks that do not interact at all. Exact compression of interacting agents can come from symmetry, as when every member of a block can be rotated together without changing any cost, or from special energy structure such as the center-of-mass cancellation; these are two sufficient routes, and the general test is whether the retained velocity is constant along each fiber. A related consequence is that exact inheritance creates no oscillation at the image of an equilibrium of the fine flow: the linearization of a gradient flow at a critical point has real eigenvalues, and an exact parent inherits a subset of them. A parent equilibrium need not be such an image. If one agent drifts steadily while never reaching a critical point, and the parent reads a second, relaxing agent in a frame turned by the first, the parent spirals into its equilibrium although every fine motion is a gradient descent.

\section{A missing role can prevent closure}

Let two independent represented bits have laws
$\operatorname{Ber}(p_1)$ and $\operatorname{Ber}(p_2)$. Their probabilities
evolve according to
\begin{equation}
 \begin{aligned}
 \dot p_1&=-p_1+\lambda(p_2-p_1),\\
 \dot p_2&=-2p_2+\lambda(p_1-p_2),\qquad \lambda\geq0.
 \end{aligned}
\end{equation}
The square $[0,1]^2$ is preserved: the vector field never points outward
on its boundary. It is tangent to a face when $\lambda=0$ and vanishes at the
corner $(0,0)$. Retain only the law of the sum of the two bits. The configurations
$(0.8,0.2)$ and $(0.2,0.8)$ both give probabilities
$(0.16,0.68,0.16)$ for sums $0,1,2$.

Their subsequent changes differ. The coupling terms cancel when we add the
two equations, giving
\begin{equation}
 \frac{d}{dt}\E(K_1+K_2)=-(p_1+2p_2).
\end{equation}
The derivatives are $-1.2$ and $-1.8$. Thus the same retained law has two
different possible derivatives. Compression erased which probability belonged
to the faster-decaying role. The fine recognition law remains a product
throughout. What must be retained is the information this evolution uses;
cross-agent covariance is not automatically required.

\section{An exact supplied hierarchy}

A finite construction shows that two successive reductions can work. Take
$x=(a,b,c,r)\in\{0,1\}^4$, with internal register $r$ and readout
$\rho(x)=r$. Let the agent laws be
\begin{equation}
 q_a=\operatorname{Ber}\left(\frac{1+2a}{4}\right),\qquad
 s_a=\operatorname{Ber}\left(\frac{2+a}{4}\right).
\end{equation}
The model evaluator assigns probability $1/4+m/2$ to record value one
when the model presentation is $m\in\{0,1\}$. Averaging over $s_a$
gives the predictive law $\nu_a=\operatorname{Ber}(1/2+a/8)$.

At each step draw independent bits $\xi,\eta,\upsilon$ with probabilities
$1/4,1/3,1/5$ of being one, independently of previous draws, and set
\begin{equation}
 \begin{aligned}
 a'&=a\mathbin\oplus\xi,\\
 b'&=b\mathbin\oplus a\mathbin\oplus\eta,\\
 c'&=c\mathbin\oplus b\mathbin\oplus\upsilon.
 \end{aligned}
\end{equation}
Here $\oplus$ means addition modulo two. Independently of these draws,
conditional on the current state, give $r'$ law $\nu_a$. The record is
therefore an internal state change. Its generation rule is supplied, rather
than deduced from the existence of an evaluator.

First retain $(a,b,r)$, then retain $(a,r)$. Summing the fine transition
over $c'$ removes a normalized factor and leaves a row independent of $c$.
Summing that row over $b'$ removes another normalized factor and leaves a
row independent of $b$. The configuration spaces shrink from sixteen to
eight to four states. Both parents retain $q_a,s_a$, their evaluator, and
the register. They reproduce the complete retained state and record
histories for every initial law. The coupling is real, since $a$ affects
$b'$ and $b$ affects $c'$; its triangular form permits these reductions.
This is a drive and its response: $a$ drives $b$, which drives $c$, and
nothing acts back. Each retained block is autonomous because the variables
it keeps never read the variables it discards. The configuration describes
one agent with internal marks, so neither level groups two agents, and the
hierarchy itself was prescribed.

\section{What a parent must remember}

When closure fails, a parent can retain more state or, under a declared
Markov process, the conditional law of the current hidden state given its
accessible history. The Markov property matters: if the next state can depend
on older states directly, two histories with the same current conditional law
can predict different futures. In a finite model this is a probability
vector: propagate it with
the fine transition, restrict to the new retained value, and normalize.
That vector determines every future response even when the current summary
does not. Keeping only an equilibrium hidden distribution generally changes
transient predictions. The vector lives on the simplex of all fine states,
so it saves nothing unless noise or poor observability keeps the history from
pinning down the fine state.

There is a middle case in which the parent needs no memory although the
current summary is not autonomous for every preparation. Suppose that for
each retained value $z$ there is a fixed way $\Lambda(z,\cdot)$ of spreading
probability over the fine states with that value, and that spreading and then
evolving gives the same result as evolving the parent and then spreading,
$\Lambda T=T^c\Lambda$. A child population prepared by spreading a parent law
in this way stays spread in this way forever: the hidden distribution is
always $\Lambda(Z_n,\cdot)$, and the parent is an ordinary Markov chain with
kernel $T^c$ \citep{RogersPitman1981}. A handoff that prepares children
differently can break the parent, even though the kernels are unchanged.

Given the fine update and a list of what the parent must reproduce (its
laws, the responses to a declared set of probes, and the channels it reads),
there is a smallest exact parent. For deterministic dynamics two
configurations are lumped together exactly when the listed quantities agree
for them now and at every later time. Nothing coarser can be exact, and
anything exact remembers at least this much. For random dynamics the
smallest exact parent is finer: two states are lumped when they send equal
probability, jointly with each listed response, into every lumped class,
which can separate states whose future listed quantities have the same
laws. A longer list of probes gives a finer parent; if outside agents
may probe the children without restriction, they can tell every exposed child
law apart and the parent cannot forget any of them. How much a parent
forgets is therefore decided by the probes declared before the parent is
chosen. ``Before'' has a definite meaning here: the probes are given by a
rule that assigns a class to every possible block from the constitutive data
alone. A list written down after the retention, such as ``every probe whose
answer the parent can reproduce'', would make any retention exact.

There is an exact way to separate insufficient memory from inaccurate prediction. Let $P$ be a declared law of a finite record sequence $Z_0,\ldots,Z_{T-1}$ on standard Borel spaces. A retained statistic is $M_t=f_t(Z_{<t})$. Let $S$ use normalized prediction rows $s_t(\cdot\mid M_t)$. The path-law chain rule and conditional disintegration give
\begin{equation}
\begin{split}
\KL(P\|S)
={}&\sum_{t=0}^{T-1}I_P(Z_t;Z_{<t}\mid M_t)\\
&+\sum_{t=0}^{T-1}\E_P
\KL\bigl(P(Z_t\in\cdot\mid M_t)\|s_t(\cdot\mid M_t)\bigr).
\end{split}
\end{equation}
The first sum is information about the next record that was discarded with the history. The second is error in the chosen prediction given the retained information. Both are nonnegative, including when the total is infinite. This extends the finite-window Markov decomposition \citep[Thm.~7.3.1]{Gray1990}. Zero first sum gives predictive sufficiency for the declared process; a compact recursively updated memory still needs a realization $M_{t+1}=U(M_t,Z_t)$. Changing the allowed controls changes which histories must be predicted.

Statistical dependence alone also does not select an agent hierarchy. A declared finite cross-agent joint law has a finest partition into independent blocks, but arbitrarily weak couplings can merge those blocks. A deterministic selector reading only a fully permutation-symmetric law cannot choose a nontrivial intermediate hard partition while respecting relabeling: swapping one member with an outside agent would create a different block that overlaps the original. More generally, such a selector can only output a partition that every symmetry of the law preserves; on a ring of four agents with the symmetries of a square, it can pair opposite agents but never neighbors. Realized histories, additional geometry, or a distribution over partitions can fall outside that restriction. The result identifies what a selector must read; it does not identify independence with agency.

\section{The formation standard}

The theory aims to show that interacting agents can produce further agents.
A construction that groups agents and hands their laws to parents does not
yet show this, because the grouping and the handoff can be arranged without
the interaction doing anything. The formation standard of the main text
lists six clauses, each closing one such loophole, and it asks that every
choice that could otherwise be made after seeing the dynamics be fixed
first: the controls, the sense of ``open'', the probes, the selector and its
constants, the time window, and the counterfactual handoffs.

The first clause asks for a set of starting configurations that is open in
the coordinates in which the dynamics is actually posed, or open inside a
family preserved by the dynamics and defined by stated constraints, with
that weaker grade reported. A single trajectory is not enough. Along one
path, almost any summary of the state traces out a curve, and every curve is
the solution of some differential equation, so ``the parent evolves
autonomously'' would hold for nearly any choice of parent. Exactness means
something only when it is tested on configurations that the parent cannot
tell apart. Declaring an arbitrary topology would not help, since a single
point is open in the discrete topology. The second clause asks the dynamics
to pick nontrivial groups that persist by a stated criterion, with a
selector drawn from a family whose constants are fixed in advance and range
over an interval rather than a single tuned value. The third asks that each
group carry inherited laws, evaluators, responses, and dynamics for the
probes that outsiders may use, given by the rule described in the previous
section.

The fourth clause is compression. A parent that simply keeps every child's
state is always exact and says nothing. Four agents with Gaussian laws pulled
together by symmetric peer costs make the point: a threshold rule first pairs
the two closest agents, then the other two, then the two pairs, and keeping
the full tuple of child states at each stage reproduces the dynamics exactly.
Nothing has been forgotten, so nothing has formed. The clause therefore asks
each parent to lump together configurations that a child-level probe could
still distinguish. These probes are what the agents, their channels and
records, and any declared external querier can do to individual children
before the handoff. A distinction that none of them can resolve is gauge in
the sense of the first chapter, and erasing it is not compression.

The fifth clause is the control experiment. The main text declares two
nulls. In the label-preserving null every peer cost is deleted while the
sources keep their labels and prior weights; a receiver that also compares
itself with its own model then gives that self comparison a weight below
one, because the deleted peers still hold part of its row at zero cost. In
the renormalized null the deleted peers are removed from the row, so the self
comparison keeps weight one. When an agent's self cost is nonzero the two
nulls change the self force differently, and a grouping difference that
appears against only one of them may come from that reweighting rather than
from the interaction. The standard therefore asks for a difference against
both. Only costs that actually read another agent's law are deleted. A cost
whose comparison law ignores every other agent, such as one through a
channel with a fixed output, stays in place: deleting it would move each
agent's own target, and agents that never interact could then pass the
control. A third control, the frozen environment, deletes nothing but holds
every other agent's law at its starting value in each agent's force. It asks
a different question, whether each agent's response to where the others
started already produces the groups, without any feedback.

The standard asks that the sequence of selected groups differ from the null
one at every point of an open set of starting configurations. Only the
groups count. If the interaction merely makes the same joins happen sooner
or later, or reverses the order of two joins of disjoint groups, it has
changed a clock, not the outcome. On a finite time window the comparison has
to be made with care: if the interaction doubles the speed of a common
clock, the full flow can complete within the window joins that the null
completes only later. The standard therefore compares the full sequence on
its window with the beginning of the complete null sequence. Agents that do
not interact can be grouped by any rule that reads their initial data and a
common relaxation clock, much as runners in a race can be grouped by
finishing order even though they never touch, and the standard excludes
exactly that.

The verdict of the fifth clause is reported with the selector constants for
which it holds and with a quantity that does not depend on the selector at
all. Without that discipline a threshold placed after the trajectories are
known can decide the verdict. A stronger grade asks for a difference at every
starting point and for every constant, and an amplifying grade asks that
groups be tighter relative to their separations than the same agents, on
the same partition, at the same times, without interaction.

A whole class of architectures fails this control. If every agent has the
same model-independent evaluator, no agent attends to itself, and every
channel delivers one common law that depends on the population only through
an additive moment, then each pair of agents relaxes toward each other at a
common rate: the logarithm of the ratio of their laws shrinks by the factor
$e^{-2t}$, up to a constant. Without the peer costs the same quantity shrinks
by $e^{-t}$. More precisely, the full flow is the null flow run at doubled
speed and then tilted by one common factor. Any selector that is blind to
such a common tilt, for example one that compares pairs through their law
ratios, produces the same sequence of groups with or without the interaction,
at different times. A selector that is not blind to it, for example a
threshold on the divergence between two laws, can give different groups,
which is why the selector is fixed in advance and a selector-free quantity
is reported. Once agents also compare themselves with their own models
through the common channel, the two nulls no longer share a clock, and the
argument does not apply. A related statement holds for scalar parameters
moved by a common state: if summing two agents' parameters is an exact parent
for every pair, the parameter equation is affine, all differences shrink by
a common factor, and the order of joins along the full flow is fixed by the
initial distances. That alone does not decide the control, which also
depends on how the null flow moves.

The sixth clause is recursion. The parents must themselves be grouped, by
the inherited interactions and the same selection rule, without resetting
their data or adding a new attraction tuned for the purpose, and this second
grouping must depend on what the first stage produced. The test uses a set
of counterfactual handoffs fixed in advance: the actual handoffs, together
with handoffs made at the same times from runs in which the first groups had
not yet formed. If the second grouping happens from all of them except a set too thin to contain an open piece of the class, or from every one of
them, it is a sequence of scales supplied in the initial data, not a parent
built from parents. Counterfactual parents that already sit inside the
joining threshold are left out of this set, since a rule that joins groups
when their distance falls through a threshold never joins them, and
including them would let every such construction pass.

A construction that meets the selection, inheritance, and first recursion
requirement but is shown to fail compression, the control, or the
counterfactual test is called a two-stage selection with exact inheritance.
When a clause is neither established nor shown to fail, its status is left
open and named. The standard does not ask that a parent be an agent of the
same type as its children or that it act back on them, so two stages here do
not yet mean a hierarchy that repeats itself.

The main text calibrates the clauses on four Gaussian agents pulled toward
their own fixed anchors and toward each other. The coupling shrinks every
difference to a third of the corresponding anchor difference, so two pairs
form and then the pair of pairs, while without the coupling nothing forms at
all; the control is passed, even in its strong grade. The frozen
environment already reproduces the groups. And the union of the two parents
forms from every counterfactual handoff, because the gap between two parents
always relaxes to the same value set by the anchors. The example is
therefore a two-stage selection that fails the counterfactual test, which is
what it should be: its stages are supplied by the anchors.

The constructions in the later chapters are measured against this standard,
each with a report of the same form. The nonsmooth fusion example produces
identical fused children, so it compresses nothing. The common-precision
matrix-law construction of \cref{cf:formation} selects two parents and then
their union with a schedule that does not depend on the coupling; there the
two nulls differ, because each agent compares itself with its own model, and
the control passes against one and fails against the other, so it fails. The
six-agent construction of \cref{cf:open-formation} works on a full open
neighborhood of laws with selective channels, but its theorem does not
settle the control or the counterfactual test, and integrations of a limit
field indicate that the null flow forms the same groups. Both smooth
constructions supply their compression architectures in advance. No
construction in this manuscript is proved to meet every clause.

% !TeX root = ../MAgent_Theory_2.0_companion.tex
\chapter{Links and Physical Geometry}
\label{comp:geometry}

The agent graph and the context base describe different directions of comparison. An edge $j\to i$ compares agents at one context. A path in the base $\mathcal C$ compares states at successive contexts. A product of communication links can be nontrivial without defining curvature on $\mathcal C$. Conversely, one agent can carry a curved connection over contexts; collective averaging is not required.

\section{Relating neighboring contexts}

To construct that connection, suppress the sector label and fix one agent. On a coordinate neighborhood, take a twice continuously differentiable presentation $U(c)$ and a continuously differentiable matrix-valued one-form $a=a_\mu(c)dc^\mu$. Declare directional links whose first-order behavior is
\begin{equation}
\Theta(c+\delta c,c)
=U(c+\delta c)e^{-a_\mu(c)\delta c^\mu}U(c)^{-1}
+O(|\delta c|^2).
\end{equation}
The displacement enters linearly at first order, and zero displacement gives identity. These are additional properties beyond having a finite interaction matrix at each context. For a principal $G$-connection, the generators must belong to the represented Lie algebra; unrestricted matrices instead define general linear bundle connections.

Using $D=d+B$, expansion identifies
\begin{equation}
B=UaU^{-1}-dUU^{-1}.
\end{equation}
The first term carries the declared directional interaction into displayed coordinates. The second compensates for the change of presentation along the path. Their relationship is exactly the operator identity
\begin{equation}
d+B=U(d+a)U^{-1}.
\end{equation}
Indeed, acting on a vector field $f$ gives $U d(U^{-1}f)=df-dUU^{-1}f$; the remaining term is $UaU^{-1}f$. Squaring the operator therefore gives
\begin{equation}
F=dB+B\wedge B=U(da+a\wedge a)U^{-1}.
\end{equation}
This is the mathematical connection to gauge theory, with the sign convention fixed by the displayed covariant derivative \citep{Tong2018}.

\section{Curvature and paths}

Curvature need not come from spatially varying coefficients. On a two-dimensional context patch choose $U=I$, $a_x=E_{12}$, and $a_y=E_{21}$. Both directional generators have rank one, but
\begin{equation}
F_{xy}=[E_{12},E_{21}]
=\begin{pmatrix}1&0\\0&-1\end{pmatrix}.
\end{equation}
The matrix $E_{ab}$ has one unit entry in position $(a,b)$. Applying the two infinitesimal operations in opposite orders gives different results. This example supplies directional data; it does not derive them from the agents' laws. A presentation change alone, $B=-dUU^{-1}$, instead has zero curvature.

Finite path transport solves
\begin{equation}
\frac{dT}{dt}=-B(\dot\gamma(t))T(t),\qquad T(0)=I.
\end{equation}
Its solution is a path-ordered exponential. Concatenating paths composes their transports; reversing the same path gives the inverse. Two independently specified attending edges can have unrelated generators, so their opposite labels do not impose this inverse relation. The distinction is necessary when borrowing Wilson-loop language from gauge theory \citep{Wilson1974}.

\section{Local descriptions and global compatibility}

Endpoint covariance requires
\begin{equation}
\Theta'(c_2,c_1)=g(c_2)\Theta(c_2,c_1)g(c_1)^{-1}.
\end{equation}
Differentiating the endpoints produces the affine law
\begin{equation}
B'=gBg^{-1}-dg g^{-1}.
\end{equation}
Likewise, keeping the intrinsic transport fixed under $U'=Uh$ requires $a'=h^{-1}ah+h^{-1}dh$. A directional receiving--sending product that transforms only by conjugation is a tensorial one-form. It can supply the difference between two connections, but does not by itself supply the full affine transformation law.

Global compatibility adds another requirement. If patch coordinates obey $v_\alpha=g_{\alpha\beta}v_\beta$, then $g_{\alpha\beta}g_{\beta\gamma}=g_{\alpha\gamma}$ on triple overlaps, and local connection forms must satisfy the affine law there. A nontrivial bundle may have no single global presentation. Even a trivial bundle can carry flat transport with global memory: on a circle, the scalar connection $B=b d\theta$ has $F=0$ but holonomy $e^{-2\pi b}$ around one turn. For real $b\ne0$, this differs from identity.

Physical interpretation requires a model of the contexts, an agent mechanism producing compatible directional links, and observable responses to their transport. A connection alone does not supply spacetime or quantum dynamics.

\appendix
% !TeX root = ../MAgent_Theory_2.0_companion.tex
\chapter{Notation in Physical Terms}
\begingroup
\small\setstretch{0.94}
\renewcommand{\arraystretch}{1.0}
\begin{longtable}{@{}>{\raggedright\arraybackslash}p{0.19\textwidth}>{\raggedright\arraybackslash}p{0.75\textwidth}@{}}
\toprule
Symbol & Meaning and role\\
\midrule\endhead
\bottomrule\endfoot
$k,m$ & A represented state possibility and a model presentation. Neither symbol by itself denotes a probability law.\\
$q_i,s_i$ & The agent's belief about represented contents (its perception) and its belief over generative models. They can evolve in coupled equations.\\
$\zeta_i$ & The joint law of represented state and model contents; $q_i\otimes s_i$ is a declared product restriction.\\
$g^F,\ell_v$ & Fisher distinguishability on a regular statistical family and the score describing one infinitesimal change.\\
$D^\omega\zeta,h^\omega_\zeta$ & Statistical change after contextual transport is accounted for, and the resulting tensor on the context base.\\
$A^F,\mathfrak k$ & The collectively minimizing connection class and the metric measuring common infinitesimal frame motions.\\
$G_i(m,dk)$ & An evaluator assigning state predictions to each model presentation. Averaging gives $s_iG_i$.\\
$r_i,\ell_i,P_i$ & Prior over presentations, record likelihood, and the normalized generative model $r_iG_i\ell_i$. A model law predicts records only through $G_i$ and $\ell_i$.\\
Species & Agents that share a generative model up to a change of frame. They may still hold different beliefs.\\
Three objectives & A global free energy of one normalized joint over all agents; per-agent free energies in which peers' laws are data; the reciprocal law potential, which moves sources toward their receivers and is not a free energy of any fixed model with the agent laws as recognition laws.\\
$\mathcal F_i$ & Free energy of agent $i$'s own generative model. In the per-agent picture its own evaluator and its peers' channelled laws compete as sources of prior.\\
$x\in\mathsf S$ & A complete configuration of agents, including their laws and the memory and interfaces read by the dynamics.\\
$\rho_i:\mathsf S_i\to\mathsf O_i$ & A readout of an internal distinction. Its value is a record.\\
$\Gamma_{ij}$ & A normalized interaction transition on complete configurations. It induces record probabilities through the readout.\\
$U_i^a$ & An invertible frame map for sector $a=q,s$. Changing frame changes the description.\\
$T_j^a,O_i^a$ & Sending and receiving interfaces. Their product can discard directions.\\
$L_{ij}^a,\Theta_{ij}^a$ & A reference generator and its framed exponential interaction. The generator may be singular while its exponential is invertible.\\
$\beta_{ij},\gamma_{ij}$ & Belief-sector and model-sector source weights in normalized rows.\\
$\mathcal E$ & The reciprocal law potential on agent laws. Its gradient dynamics requires a metric, and every edge pushes on both its receiver and its source.\\
$\mathcal F$ & Variational free energy for a specified normalized probabilistic model and conditioned record.\\
$Q_xC$ & A law of compressed represented contents. It does not specify a parent update.\\
$R(x),F^c$ & A complete retained state and its inherited update. They must satisfy $R\circ F=F^c\circ R$, with $\circ$ composing maps in function order.\\
$\Phi^0$ & The null flow: the same construction with every peer cost deleted. It is the control experiment for selection.\\
$P,P^0$ & Partition sequences selected along the flow and along the null flow, without event times.\\
(F5) & Interaction necessity: $P\neq P^0$ on an open set of admitted initial data. A shift of event times alone does not count.\\
Two-stage selection & Short for two-stage selection with exact inheritance: selection, inheritance, and a second selection among parents, with at least one of compression, interaction necessity, or counterfactual dependence of the second stage failing. It is not a formation, and the failed clauses are stated.\\
$W_2(\nu,\mathsf P)$ & A Wishart law of positive-definite matrix contents, with real shape $\nu>1/2$ and precision $\mathsf P$. It is not a Gaussian closure.\\
$\alpha_i,\eta_i; A_I,B_I,n_I$ & Individual belief and model shapes, their block sums, and the number of primitive children retained by a parent.\\
$B^q_{IJ},B^s_{IJ}$ & Native source masses between current blocks, averaged over receivers with primitive-count weights.\\
$\mathsf H_I,\phi_I$ & Relative constant angular offsets discarded by a block, and its retained local mean phase at the contextual anchor.\\
$\mathcal C,B,F$ & The context base, a connection one-form, and its curvature. These require a rule for relating neighboring contexts. The symbol $F$ also denotes a discrete update in the inheritance chapter.\\
\end{longtable}
\endgroup

\bibliographystyle{plainnat}
\bibliography{references}
\end{document}
